\documentclass[12pt,a4paper]{report}

%% ─── Packages ───────────────────────────────────────────────────────
\usepackage[utf8]{inputenc}
\usepackage[T1]{fontenc}
\usepackage[french]{babel}
\usepackage[top=2cm,bottom=2cm,left=2.5cm,right=2.5cm]{geometry}
\usepackage{amsmath,amssymb}
\usepackage{graphicx}
\usepackage{xcolor}
\usepackage{listings}
\usepackage{fancyhdr}
\usepackage{booktabs}
\usepackage{tcolorbox}
\usepackage{hyperref}

%% ─── Couleurs ───────────────────────────────────────────────────────
\definecolor{Navy}{RGB}{10,36,99}
\definecolor{Steel}{RGB}{70,130,180}
\definecolor{DGreen}{RGB}{0,110,55}
\definecolor{CodeBg}{RGB}{248,249,252}

%% ─── Code listings ──────────────────────────────────────────────────
\lstdefinestyle{mystyle}{
  language=Python,
  backgroundcolor=\color{CodeBg},
  basicstyle=\small\ttfamily,
  keywordstyle=\color{Navy}\bfseries,
  stringstyle=\color{DGreen},
  commentstyle=\color{gray}\itshape,
  breaklines=true,
  tabsize=2,
  numbers=left,
  numberstyle=\tiny,
}
\lstset{style=mystyle}

%% ─── Headers ────────────────────────────────────────────────────────
\pagestyle{fancy}
\fancyhf{}
\fancyhead[L]{\small\textbf{LightGCN - Recommandation}}
\fancyhead[R]{\small Kenné Takoudjou — 24P817}
\fancyfoot[C]{\thepage}
\renewcommand{\headrulewidth}{0.5pt}

%% ═════════════════════════════════════════════════════════════════════
\begin{document}

%% ════════════════════════════════════════════════════════════════════
%%  PAGE DE TITRE SIMPLE
%% ════════════════════════════════════════════════════════════════════
\begin{titlepage}
\centering
\vspace*{2cm}

{\Large\bfseries\color{Navy}
  ÉCOLE NATIONALE SUPÉRIEURE POLYTECHNIQUE DE YAOUNDÉ\\
  ENSPY — Université de Yaoundé I
}

\vspace{2cm}
{\color{gray}\hrule}

\vspace{2cm}
{\Huge\bfseries\color{Navy}
  Système de Recommandation\\
  avec LightGCN
}

\vspace{0.5cm}
{\Large\color{Steel}
  Graph Convolutional Networks pour MovieLens
}

\vspace{2cm}
{\color{gray}\hrule}

\vspace{3cm}

\begin{tabular}{ll}
  \textbf{Étudiant :}     & Kenné Takoudjou Franck Manuel \\
  \textbf{Matricule :}    & 24P817 \\
  \textbf{Classe :}       & AIA4 — Art et Intelligence Artificielle \\
  \textbf{Professeur :}   & M. Bitha \\
  \textbf{Année :}        & 2025/2026 \\
\end{tabular}

\vfill
{\small\color{gray}\today}

\end{titlepage}

%% ════════════════════════════════════════════════════════════════════
%%  RÉSUMÉ EXÉCUTIF
%% ════════════════════════════════════════════════════════════════════
\newpage
\section*{Résumé exécutif}
\addcontentsline{toc}{section}{Résumé exécutif}

Ce projet implémente un système complet de recommandation de films basé sur
\textbf{LightGCN} (Light Graph Convolutional Networks), un modèle de deep
learning pour graphes bipartis. Le système utilise le dataset \textbf{MovieLens
small} contenant 610 utilisateurs et 9 724 films.

\begin{center}
\begin{tcolorbox}[colback=yellow!10,colframe=Navy,arc=3pt,boxrule=1pt]
\textbf{Résultats clés :}
\begin{itemize}
  \item Recall@10 = \textbf{0.1187}
  \item NDCG@20 = \textbf{0.1943}
  \item 510 208 paramètres apprenables
  \item Temps d'entraînement : 572 secondes (30 époques)
\end{itemize}
\end{tcolorbox}
\end{center}

\tableofcontents

%% ════════════════════════════════════════════════════════════════════
%%  CHAPITRE 1 — INTRODUCTION
%% ════════════════════════════════════════════════════════════════════
\chapter{Introduction}

\section{Motivation}

Les systèmes de recommandation sont essentiels pour les plateformes en ligne
(streaming, e-commerce, réseaux sociaux). Deux approches dominent :

\begin{itemize}
  \item \textbf{Filtrage basé contenu} : recommande des items similaires aux
        préférences passées
  \item \textbf{Filtrage collaboratif} : exploite les interactions collectives
        pour découvrir des patterns cachés
\end{itemize}

LightGCN combine le filtrage collaboratif avec les réseaux de neurones
convolutifs sur graphes, offrant une approche élégante et efficace.

\section{Objectifs du projet}

\begin{enumerate}
  \item Construire un graphe biparti utilisateur-film
  \item Implémenter LightGCN avec propagation multicouches
  \item Entraîner le modèle avec la perte BPR
  \item Évaluer les performances (Recall@K, NDCG@K)
  \item Déployer une API REST avec Flask
\end{enumerate}

%% ════════════════════════════════════════════════════════════════════
%%  CHAPITRE 2 — FONDEMENTS THÉORIQUES
%% ════════════════════════════════════════════════════════════════════
\chapter{Fondements théoriques}

\section{Architecture LightGCN}

\subsection{Graphe biparti utilisateur-item}

Les interactions forment un graphe $\mathcal{G} = (\mathcal{U} \cup \mathcal{I},
\mathcal{E})$ où :
\begin{itemize}
  \item $\mathcal{U}$ : ensemble des utilisateurs ($N = 609$)
  \item $\mathcal{I}$ : ensemble des films ($M = 7\,363$)
  \item $(u,i) \in \mathcal{E}$ : si l'utilisateur $u$ a noté le film $i \geq 3.5$
\end{itemize}

\subsection{Matrice d'adjacence normalisée}

La matrice d'adjacence bipartite est :
\begin{equation}
  \mathbf{A} = \begin{pmatrix} \mathbf{0}_{N\times N} & \mathbf{R} \\
                               \mathbf{R}^\top & \mathbf{0}_{M\times M}
              \end{pmatrix}
\end{equation}

Après normalisation symétrique : $\tilde{\mathbf{A}} = \mathbf{D}^{-1/2}
\mathbf{A} \mathbf{D}^{-1/2}$ où $\mathbf{D}$ est la matrice des degrés.

\subsection{Propagation de messages}

À chaque couche $k$, les embeddings se propagent via :
\begin{equation}
  \mathbf{E}^{(k+1)} = \tilde{\mathbf{A}} \cdot \mathbf{E}^{(k)}
\end{equation}

Les embeddings finals sont obtenus par agrégation moyenne sur $K+1$ couches :
\begin{equation}
  \mathbf{e}^* = \frac{1}{K+1}\sum_{k=0}^{K}\mathbf{e}^{(k)}
\end{equation}

\subsection{Score de prédiction et perte BPR}

Score d'affinité :
\begin{equation}
  \hat{y}_{ui} = {\mathbf{e}_u^*}^\top \mathbf{e}_i^*
\end{equation}

Perte d'entraînement (BPR) :
\begin{equation}
  \mathcal{L} = -\sum \ln\sigma(\hat{y}_{ui^+} - \hat{y}_{ui^-}) + \lambda\|\Theta\|_2^2
\end{equation}

\subsection{Avantages de LightGCN}

\begin{center}
\begin{tabular}{lcc}
\toprule
\textbf{Caractéristique} & \textbf{GCN} & \textbf{LightGCN} \\
\midrule
Transformations linéaires & ✓ & ✗ \\
Fonctions d'activation & ✓ & ✗ \\
Self-connections & ✓ & ✗ \\
Nombre de paramètres & $O(N \cdot d \cdot K)$ & $O((N+M) \cdot d)$ \\
\bottomrule
\end{tabular}
\end{center}

LightGCN supprime les éléments non essentiels, réduisant drastiquement les
paramètres tout en améliorant les performances.

%% ════════════════════════════════════════════════════════════════════
%%  CHAPITRE 3 — MISE EN ŒUVRE
%% ════════════════════════════════════════════════════════════════════
\chapter{Mise en \oe{}uvre}

\section{Préparation des données}

\subsection{Dataset MovieLens small}

\begin{center}
\begin{tabular}{ll}
\toprule
\textbf{Métrique} & \textbf{Valeur} \\
\midrule
Ratings total & 100\,836 \\
Ratings filtrés ($\geq 3.5$) & 61\,716 \\
Utilisateurs & 609 \\
Films & 7\,363 \\
Densité & 1.38\% \\
\bottomrule
\end{tabular}
\end{center}

\subsection{Splits train/test}

Split stratifié par utilisateur : 80\% train / 20\% test.

\begin{itemize}
  \item \textbf{Train} : 49\,619 interactions
  \item \textbf{Test} : 12\,097 interactions
  \item \textbf{Utilisateurs avec test} : 604
\end{itemize}

\section{Architecture du modèle}

\begin{lstlisting}[caption=Classe LightGCN simplifié]
class LightGCN(nn.Module):
  def __init__(self, n_users, n_items, embed_dim, n_layers, adj):
    super().__init__()
    self.adj = adj
    self.n_layers = n_layers
    self.user_embed = nn.Embedding(n_users, embed_dim)
    self.item_embed = nn.Embedding(n_items, embed_dim)
    nn.init.xavier_uniform_(self.user_embed.weight)
    nn.init.xavier_uniform_(self.item_embed.weight)

  def propagate(self):
    E = torch.cat([self.user_embed.weight,
                   self.item_embed.weight], dim=0)
    embs = [E]
    for _ in range(self.n_layers):
      E = torch.sparse.mm(self.adj, E)
      embs.append(E)
    return torch.stack(embs, dim=1).mean(dim=1)

  def recommend(self, users):
    all_e = self.propagate()
    return all_e[users] @ all_e[self.n_users:].T
\end{lstlisting}

\subsection{Hyperparamètres}

\begin{center}
\begin{tabular}{lll}
\toprule
\textbf{Hyperparamètre} & \textbf{Valeur} & \textbf{Raison} \\
\midrule
Dimension $d$ & 64 & Standard pour MovieLens \\
Couches $K$ & 3 & Optimum empirique (He et al., 2020) \\
Taux d'apprentissage & $10^{-3}$ & Convergence stable \\
Régularisation $\lambda$ & $10^{-4}$ & Anti-overfitting \\
Batch size & 1\,024 & Vitesse/stabilité \\
Époques & 30 & Convergence observée \\
\bottomrule
\end{tabular}
\end{center}

\section{Entraînement}

La boucle d'entraînement génère des triplets $(u, i^+, i^-)$ et optimise la
perte BPR via Adam.

\begin{center}
\begin{tabular}{cccc}
\toprule
\textbf{Époque} & \textbf{Loss} & \textbf{Recall@20} & \textbf{NDCG@20} \\
\midrule
1 & 0.2420 & 0.1551 & 0.1721 \\
10 & 0.2104 & 0.1714 & 0.1858 \\
20 & 0.1883 & 0.1711 & 0.1854 \\
\cellcolor{yellow!30}28 & \cellcolor{yellow!30}0.1701 &
  \cellcolor{yellow!30}0.1849 & \cellcolor{yellow!30}0.1990 \\
30 & 0.1671 & 0.1735 & 0.1906 \\
\bottomrule
\end{tabular}
\end{center}

Le meilleur modèle est sauvegardé à l'époque 28.

%% ════════════════════════════════════════════════════════════════════
%%  CHAPITRE 4 — RÉSULTATS ET ÉVALUATION
%% ════════════════════════════════════════════════════════════════════
\chapter{Résultats et évaluation}

\section{Métriques}

\subsection{Recall@K}

Fraction des items pertinents retrouvés dans les Top-K recommandations :
\begin{equation}
  \text{Recall@K} = \frac{|\text{Top-K recommandés} \cap \text{Items test}|}{|\text{Items test}|}
\end{equation}

\subsection{NDCG@K}

Cumulative gain normalisé (pénalise les mauvais classements) :
\begin{equation}
  \text{NDCG@K} = \frac{1}{\text{IDCG@K}} \sum_{k=1}^{K} \frac{\mathbf{1}[\text{item}_k \in \text{test}]}{\log_2(k+1)}
\end{equation}

\section{Performances finales}

\begin{center}
\begin{tcolorbox}[colback=green!10,colframe=DGreen,arc=3pt,boxrule=1pt]
\textbf{Modèle : LightGCN (K=3, d=64, 609 users, 7363 items)}

\vspace{0.3cm}

\begin{tabular}{lll}
\toprule
\textbf{K} & \textbf{Recall@K} & \textbf{NDCG@K} \\
\midrule
5 & 0.0728 & 0.1926 \\
10 & 0.1187 & 0.1879 \\
20 & 0.1793 & 0.1943 \\
50 & 0.3086 & 0.2244 \\
\bottomrule
\end{tabular}
\end{tcolorbox}
\end{center}

\textbf{Interprétation} : Le modèle retrouve environ 12\% des films pertinents
dans les 10 premières recommandations. Les films les plus pertinents sont
classés en tête (NDCG élevé).

\section{Exemples de recommandations}

\textbf{Utilisateur 0} (historique : Toy Story, Heat, Seven)

\begin{itemize}
  \item 1. Shawshank Redemption (score: 11.81)
  \item 2. Pulp Fiction (score: 11.42)
  \item 3. Forrest Gump (score: 10.73)
  \item 4. Terminator 2 (score: 8.88)
  \item 5. Twelve Monkeys (score: 8.80)
\end{itemize}

\textbf{Utilisateur 1} (historique : Shutter Island, Fight Club, Ex Machina)

\begin{itemize}
  \item 1. The Matrix (score: 4.53)
  \item 2. Lord of the Rings (score: 4.07)
  \item 3. Fight Club (score: 4.02)
  \item 4. Forrest Gump (score: 4.02)
  \item 5. Ex Machina (score: 3.91)
\end{itemize}

Les recommandations montrent une \textbf{cohérence sémantique} : le modèle
suggère des films du même genre/époque que l'historique.

%% ════════════════════════════════════════════════════════════════════
%%  CHAPITRE 5 — DÉPLOIEMENT FLASK
%% ════════════════════════════════════════════════════════════════════
\chapter{Déploiement Flask}

\section{Architecture}

L'application Flask expose l'API REST via plusieurs endpoints :

\begin{center}
\begin{tabular}{lll}
\toprule
\textbf{Route} & \textbf{Méthode} & \textbf{Description} \\
\midrule
\texttt{/} & GET & Interface web \\
\texttt{/api/recommend} & GET/POST & Recommandations Top-K \\
\texttt{/api/stats} & GET & Infos modèle \\
\texttt{/api/users} & GET & Liste utilisateurs \\
\bottomrule
\end{tabular}
\end{center}

\section{Exemple requête/réponse}

\textbf{Requête :} \texttt{GET /api/recommend?user\_id=0\&k=5}

\textbf{Réponse JSON} (extrait) :

\begin{lstlisting}[language=json]
{
  "success": true,
  "user_id": 0,
  "recommendations": [
    {
      "rank": 1,
      "title": "Shawshank Redemption (1994)",
      "score": 11.808
    },
    ...
  ],
  "model_info": {
    "parameters": 510208,
    "layers": 3,
    "n_users": 609,
    "n_items": 7363
  }
}
\end{lstlisting}

\section{Améliorations effectuées}

\begin{itemize}
  \item Conversion robuste de types NumPy → Python (JSON-sérialisable)
  \item Gestion complète des erreurs avec logging
  \item Validation des entrées utilisateur
  \item Réponses structurées avec flag \texttt{success}
  \item Support GET et POST
\end{itemize}

%% ════════════════════════════════════════════════════════════════════
%%  CHAPITRE 6 — CONCLUSION
%% ════════════════════════════════════════════════════════════════════
\chapter{Conclusion}

\section{Synthèse}

Ce projet a implémenté avec succès un système complet de recommandation basé
sur LightGCN :

\begin{enumerate}
  \item \textbf{Modèle} : 510K paramètres, 3 couches, dimension 64
  \item \textbf{Performance} : Recall@10 = 11.87\%, NDCG@20 = 19.43\%
  \item \textbf{Déploiement} : API REST Flask fonctionnelle
  \item \textbf{Qualité} : Recommandations sémantiquement cohérentes
\end{enumerate}

\section{Résultats clés}

\begin{center}
\begin{tcolorbox}[colback=blue!5,colframe=Navy,arc=3pt,boxrule=1pt]
\begin{tabular}{ll}
  \textbf{Architecture} & LightGCN (K=3, d=64) \\
  \textbf{Dataset} & MovieLens small (609 users, 7,363 items) \\
  \textbf{Interactions train} & 49,619 \\
  \textbf{Meilleur Recall@20} & 0.1849 (époque 28) \\
  \textbf{Temps d'entraînement} & 572 secondes (CPU) \\
\end{tabular}
\end{tcolorbox}
\end{center}

\section{Perspectives futures}

\begin{itemize}
  \item Augmenter la dimensionalité d'embedding ($d = 128$)
  \item Expérimenter avec plus de couches ($K = 4, 5$)
  \item Intégrer des features de contenu (genres, année)
  \item Déployer avec GPU pour améliorer la vitesse
  \item Ajouter un frontend interactif
\end{itemize}

%% ════════════════════════════════════════════════════════════════════
\end{document}
